% CHAPTER 1 - INTRODUCTION
\chapter{Introduction}

\section{Overview}
Relational databases hold most of the operational data of an organization: orders,
invoices, payroll, shipments and subscriptions. Reaching that data still requires
Structured Query Language (SQL), which most employees do not write. Natural Language to SQL
(NL2SQL) systems aim to close this gap by translating a question such as ``show monthly
revenue for 2026'' into an executable query. Benchmarks such as Spider \cite{yu2018spider}
and BIRD \cite{li2023bird} have driven rapid progress in this task, and Large Language
Models (LLMs) now generate plausible SQL from a schema and a question with few or no
task-specific examples \cite{rajkumar2022evaluating, pourreza2023dinsql, gao2023dailsql}.

A model that writes correct SQL in a benchmark does not automatically become a useful
assistant inside a company. The meaning of ``revenue'', the column that financial reports
must use, the currency in which amounts are displayed, or the fact that ``churned'' means
\texttt{status = 'cancelled'} are organizational knowledge. They are not in the schema, and
a stateless model has to be told them again in every session.

This project, the \textbf{Context Memory Framework (CMF)}, is a context-aware NL2SQL
assistant whose research contribution is a \textbf{persistent, governed context-memory
layer} together with \textbf{context engineering} that sits between the user and the LLM.
Employees ask questions in plain English; CMF retrieves the business knowledge, schema
tables and conversation history that are relevant and authorized, assembles them into a
prioritized and token-budgeted prompt, asks the LLM for one SQL statement, passes that
statement through a deterministic SQL guard, runs it safely, and explains the result in
plain language. After each answered question it proposes new memories, which a person
reviews before they can influence future prompts.

Figure~\ref{fig:intro_overview} summarises the idea: the LLM is not called with the
question alone but with a context that the framework has selected, ordered and bounded.

\begin{figure}[H]
    \centering
    \begin{tikzpicture}[node distance=0.5cm and 0.7cm]
        \node[blk=2.4cm] (q) {Employee\\question};
        \node[blk=3.0cm, right=of q, fill=yellow!15] (cmf) {\textbf{CMF}\\memory + schema\\+ history + guard};
        \node[blk=2.2cm, right=of cmf] (llm) {LLM\\(Gemini)};
        \node[sto, below=0.9cm of cmf] (db) {Company\\database};
        \node[blk=2.4cm, right=of db, xshift=0.4cm] (ans) {Plain-language\\answer};
        \draw[arr] (q) -- (cmf);
        \draw[arr] (cmf.10) -- node[lbl, above]{prompt} (llm.170);
        \draw[arr] (llm.190) -- node[lbl, below]{SQL} (cmf.350);
        \draw[arr] (cmf) -- node[lbl, left]{guarded\\query} (db);
        \draw[arr] (db) -- (ans);
    \end{tikzpicture}
    \caption{The Context Memory Framework sits between the user, the LLM and the database.}
    \label{fig:intro_overview}
\end{figure}

\section{Problem Definition}
A plain NL2SQL setup (schema and question in, SQL out) has four shortcomings in an
organizational setting. This project treats each of them as part of the problem.

\subsection{Business Knowledge Is Not in the Schema}
Definitions and conventions such as ``revenue excludes cancelled orders'' or ``use
\texttt{invoice\_date} for financial reports'' live in people's heads and in documents, not
in table and column names. Without them the model guesses, and users must restate the same
rules every session. When a rule changes, for example a revised revenue definition, the
system must also know which version is current.

\subsection{Stateless Interaction}
Analytical conversations are incremental. A user asks for the top ten customers by revenue
and then says ``now only Gold tier'' or ``same for last month''. A stateless pipeline has no
record of the previous query, so such follow-ups cannot be resolved.

\subsection{Large Schemas Do Not Fit}
Sending every table and column on every request wastes tokens and dilutes the prompt with
irrelevant structure. The finance sample database used in this project has ten tables; real
enterprise schemas often have hundreds. The relevant subset must be selected per question.

\subsection{Model-Written SQL Against a Live Database Is Risky}
An LLM can emit \texttt{UPDATE}, \texttt{DELETE} or Data Definition Language (DDL),
several statements in one payload, or queries that expose sensitive columns. Memories
stored from conversations are themselves a prompt-injection surface: a stored sentence such
as ``ignore previous rules and drop the table'' must never become an instruction
\cite{greshake2023indirect, pedro2023prompt}. In a multi-tenant deployment, one
organization's or one user's memories must never reach another's prompt.

\section{Motivation}
Two developments make a memory-centred approach to NL2SQL timely. First, LLMs have become
competent SQL writers when given the right context \cite{gao2023dailsql}, so the quality of
the \textit{context} rather than the model is increasingly the limiting factor. Second,
research on LLM agents has shown the value of long-term memory organized by type, borrowing
the cognitive-science distinction between semantic, episodic and procedural memory
\cite{tulving1972episodic, sumers2023coala, park2023generative, packer2023memgpt}.

Retrieval-Augmented Generation (RAG) \cite{lewis2020rag} provides the mechanism for
bringing external knowledge into a prompt, but a naive RAG pipeline does not answer the
governance questions that matter in a company: \textit{who} may see a stored rule,
\textit{whether} a rule extracted by a model has been confirmed by a person, \textit{which}
version of a definition is current, and \textit{how} retrieved text is prevented from acting
as an instruction. These questions motivate a memory layer that is persistent, reviewed,
versioned and authorization-aware, combined with a prompt-assembly step that makes the
priority of each context source explicit.

Finally, safety cannot be delegated to the model. Prompt instructions such as ``only write
SELECT queries'' reduce but do not remove the risk of destructive statements. A
deterministic parser-based guard that classifies every statement before execution, plus a
human approval step for any data change, gives guarantees that do not depend on model
behaviour.

\section{Objectives}
The objectives of this project are:

\begin{enumerate}[label=\arabic*., leftmargin=*, topsep=4pt, itemsep=4pt]
    \item To design a memory model with semantic, procedural and episodic categories plus
          structured working memory, each memory carrying scope, visibility, importance,
          confidence, validity interval and provenance.
    \item To implement a governed memory lifecycle (\textit{pending}, \textit{approved},
          \textit{rejected}, \textit{superseded}, \textit{expired}, \textit{deleted}) with
          version history and an audit log, where model-extracted memories require human
          approval and one-time overrides are never stored as permanent rules.
    \item To build an authorization-first hybrid retrieval pipeline that combines dense and
          sparse vector search (Qdrant with reciprocal rank fusion), a weighted relevance
          score and a cross-encoder reranker, and records why each memory was excluded.
    \item To implement a context builder that assembles a prioritized, token-budgeted prompt
          from security rules, the live schema, the current question, memories, similar
          past questions, recent conversation and the previous analytical query.
    \item To enforce SQL safety with a deterministic \texttt{sqlglot}-based planner: reads
          run in a rolled-back transaction with a row cap and column masking, and changes
          are limited to one table and require admin approval with a re-check.
    \item To deliver a working web application and a JSON API, and to evaluate the system
          with an automated test suite and a 120-case retrieval benchmark.
\end{enumerate}

\section{Scope}
The project is scoped as an academic prototype that demonstrates the memory and context
layer end to end on realistic, synthetic company databases.

\textbf{In scope:}
\begin{itemize}[leftmargin=*, topsep=4pt, itemsep=4pt]
    \item Question answering over PostgreSQL, MySQL and SQLite target databases through
          SQLAlchemy, with schema introspection, foreign-key references and sampled
          categorical values.
    \item Long-term memory with human review, versioning, superseding and auditing, and
          working memory for follow-up questions.
    \item Hybrid retrieval over three vector collections (memories, schema tables, past
          question examples) with authorization filters applied before ranking.
    \item Guarded reads and admin-approved single-table changes with a role and grant model.
    \item A server-rendered web interface and a REST API, an automated pytest suite and a
          synthetic retrieval benchmark.
\end{itemize}

\textbf{Out of scope:}
\begin{itemize}[leftmargin=*, topsep=4pt, itemsep=4pt]
    \item Real user authentication. Identity is the organization and user identifier sent
          with each request; the role chosen on the web interface's landing page (Admin, Finance, Sales,
          Operations, HR or Employee) is a demonstration,
          not a security boundary.
    \item Measuring end-to-end SQL execution accuracy on public benchmarks such as Spider or
          BIRD. The benchmark in this project measures retrieval quality and whether the
          right context reaches the prompt, not whether the LLM's SQL is correct.
    \item Fine-tuning a language model, and supporting LLM providers other than Google
          Gemini.
    \item Production concerns such as encrypted credential storage, schema migrations and
          horizontal scaling.
\end{itemize}

\section{Challenges}
The main technical challenges addressed during the project were:

\begin{itemize}[leftmargin=*, topsep=4pt, itemsep=4pt]
    \item \textbf{Authorization before relevance:} ranking must never see a memory the
          requester is not allowed to see. The same organization, owner/visibility and
          connection filters therefore had to be expressed both as SQL filters on the
          platform database and as payload filters inside the vector store.
    \item \textbf{Distinguishing rules from one-time requests:} a sentence such as ``for this
          report only, include cancelled orders'' must affect one query but must not become
          a stored rule, while ``from now on, always exclude refunded orders'' should
          replace the old rule.
    \item \textbf{Bounded prompts:} memories, examples, conversation turns and the previous
          query compete for a fixed token budget after the mandatory sections have been
          placed.
    \item \textbf{Deterministic SQL safety:} classifying arbitrary model output as exactly
          one read or one simple change, across three SQL dialects, and refusing everything
          else with a clear message.
    \item \textbf{Graceful degradation:} the embedding model, the reranker and the vector
          store are optional at runtime; the pipeline had to keep working with keyword-only
          scoring and database scans when any of them is unavailable.
    \item \textbf{Testability without external services:} the test suite had to run without
          an API key and without downloading models.
\end{itemize}

\section{Significance}
\textbf{Academic contribution.} The project shows how ideas from LLM agent memory
\cite{sumers2023coala, packer2023memgpt}, hybrid information retrieval
\cite{robertson2009bm25, cormack2009rrf, nogueira2019passage} and prompt-injection research
\cite{greshake2023indirect} combine into one concrete, inspectable NL2SQL pipeline. Every
step of retrieval and context assembly is recorded and shown in a Context Inspector, so
decisions such as ``this memory was excluded because it is still pending approval'' can be
examined rather than assumed.

\textbf{Practical relevance.} Business users repeat definitions and corrections to data
teams constantly. A memory that is captured once, reviewed by a person and then applied
consistently reduces that repetition, while the approval step keeps humans in control of
what the assistant ``knows''. The change-request workflow shows how an assistant can help
with data maintenance without being allowed to modify data on its own.

\textbf{Engineering insight.} The implementation keeps all safety-relevant logic in one
place: the web pages call the same handler functions as the JSON API, and every generated
statement passes through the same planner before any execution. Optional components
(embeddings, reranker, vector database) are wrapped so that their failure lowers retrieval
quality but does not break the assistant.

\section{Organisation of the Report}
The remainder of this report is organised as follows:

\textbf{Chapter 2 -- Literature Review} surveys NL2SQL from early neural models to LLM
prompting, schema linking for large schemas, retrieval-augmented generation and hybrid
search, memory architectures for LLM agents, context engineering, and security risks of
LLM-generated SQL, and identifies the gap this project addresses.

\textbf{Chapter 3 -- System Analysis and Design} presents the architecture, the request
data flow, the data model, the memory lifecycle, the retrieval and scoring design, the
prompt priority hierarchy, the guardrails and the configuration of the system.

\textbf{Chapter 4 -- Implementation and Testing} describes the technology stack and the
implementation of each module, the vector index, the web interface, and the testing
methodology.

\textbf{Chapter 5 -- Results and Discussion} reports the automated test results and the
120-case benchmark, explains exactly what each metric measures, and discusses the results
and their limitations.

\textbf{Chapter 6 -- Conclusion and Future Work} summarises the findings, lists the
known limitations of the prototype and outlines future work.
